\PassOptionsToPackage{dvipsnames,svgnames,x11names}{xcolor}
\documentclass[tikz,border=8pt]{standalone}
\usetikzlibrary{arrows.meta,positioning,calc,shapes,shadows,decorations.pathmorphing,shapes.geometric,backgrounds,decorations.markings,decorations.pathreplacing,decorations.text,fit,shapes.misc,shapes.arrows,matrix,bending,3d,chains,intersections,shapes.symbols,svg.path,shapes.multipart,snakes,shapes.callouts,angles,quotes}
\providecolor{gold}{RGB}{212,175,55}
\providecolor{silver}{RGB}{192,192,192}
\providecolor{bronze}{RGB}{205,127,50}
\providecolor{darkblue}{RGB}{0,0,139}
\providecolor{lightblue}{RGB}{173,216,230}
\providecolor{darkgreen}{RGB}{0,100,0}
\providecolor{lightgreen}{RGB}{144,238,144}
\providecolor{darkred}{RGB}{139,0,0}
\providecolor{navy}{RGB}{0,0,128}
\providecolor{skyblue}{RGB}{135,206,235}
\providecolor{beige}{RGB}{245,245,220}
\providecolor{cream}{RGB}{255,253,208}
\usepackage{amsmath}
\usepackage{amssymb}
\usepackage{fontawesome5}
\usetikzlibrary{fadings,shadows.blur,patterns,patterns.meta}

\providecolor{gold}{RGB}{212,175,55}
\providecolor{silver}{RGB}{192,192,192}
\providecolor{bronze}{RGB}{205,127,50}
\providecolor{darkblue}{RGB}{0,0,139}
\providecolor{lightblue}{RGB}{173,216,230}
\providecolor{darkgreen}{RGB}{0,100,0}
\providecolor{lightgreen}{RGB}{144,238,144}
\providecolor{darkred}{RGB}{139,0,0}
\providecolor{navy}{RGB}{0,0,128}
\providecolor{skyblue}{RGB}{135,206,235}
\providecolor{beige}{RGB}{245,245,220}
\providecolor{cream}{RGB}{255,253,208}
\usepackage{amsmath}
\usepackage{amssymb}
\usepackage{fontawesome5}

\usepackage{tikz}
\usepackage{xcolor}

% Custom Color Palette
\definecolor{codebg}{RGB}{30, 30, 30}
\definecolor{codefg}{RGB}{212, 212, 212}
\definecolor{str}{RGB}{206, 145, 120}
\definecolor{tag}{RGB}{78, 201, 176}
\definecolor{attr}{RGB}{156, 220, 254}
\definecolor{accentR}{RGB}{220, 53, 69}
\definecolor{accentG}{RGB}{40, 167, 69}
\definecolor{accentB}{RGB}{13, 110, 253}
\definecolor{accentP}{RGB}{111, 66, 193}
\definecolor{charcoal}{RGB}{38, 38, 38}

\begin{document}
\begin{tikzpicture}[
    x=1cm, y=1cm,
    blade/.style={
        path picture={
            \fill[#1] (path picture bounding box.north west) rectangle ([xshift=4.5pt]path picture bounding box.south west);
        }
    },
    bubble/.style={
        draw=#1!50!black!30,
        top color=white,
        bottom color=#1!5,
        rounded corners=8pt,
        inner sep=14pt,
        font=\small,
        align=left,
        text=charcoal,
        line width=0.8pt,
        blur shadow={shadow opacity=12, shadow yshift=-2pt, shadow blur radius=4pt}
    },
    codebox/.style={
        fill=codebg,
        text=codefg,
        font=\ttfamily\small,
        rounded corners=8pt,
        inner sep=16pt,
        align=left,
        line width=0.8pt,
        draw=black!70,
        blur shadow={shadow opacity=20, shadow yshift=-2pt, shadow blur radius=5pt}
    },
    procbox/.style={
        draw=#1!70!black!40,
        top color=white,
        bottom color=#1!5,
        line width=1.1pt,
        rounded corners=8pt,
        inner sep=12pt,
        font=\small,
        align=center,
        text=charcoal,
        blur shadow={shadow opacity=14, shadow yshift=-2pt, shadow blur radius=4pt}
    },
    severbox/.style={
        draw=accentP,
        top color=white,
        bottom color=accentP!8,
        line width=1.8pt,
        rounded corners=8pt,
        inner sep=12pt,
        font=\small,
        align=center,
        text=charcoal,
        blur shadow={shadow opacity=28, shadow yshift=-2pt, shadow blur radius=7pt}
    },
    datapipe/.style={
        -Stealth=3mm,
        line width=2.5pt,
        draw=#1,
        shorten >= 3pt,
        shorten <= 3pt
    },
    dashedarrow/.style={
        -Stealth=3mm,
        line width=1.5pt,
        dashed,
        color=#1,
        shorten >= 3pt,
        shorten <= 3pt
    }
]

% Background Enforcement
\begin{scope}[on background layer]
    \fill[white] (-3, 0) rectangle (48, 27);
\end{scope}

% ---------------- TITLE SECTION ----------------
\node[font=\Huge\bfseries, color=charcoal, align=center] at (22, 25.5) {Mechanics of the Preload Scanner Severing};
\node[font=\Large, color=black!60, align=center] at (22, 24.5) {How \texttt{<source media="(max-width: 600px)">} mechanically halts the browser's greedy network queue.};

% Center Structural Divider
\draw[line width=1.5pt, dashed, black!20] (22, 23.5) -- (22, 1.0);


% ---------------- LEFT COLUMN (Standard <img> Pipeline) ----------------
\node[align=center] at (11, 22.5) {
    {\LARGE\bfseries\color{charcoal} Standard \texttt{<img>} Pipeline}\\[2pt]
    {\normalsize\color{black!50} (Unconditional Greedy Fetch)}
};

% Code Box
\node[codebox, text width=9.5cm] (codeL) at (11, 18.5) {
<\textcolor{tag}{img} \textcolor{attr}{src}=\textcolor{str}{"wide.jpg"}>\\
\textcolor{black!50}{<!-- No media conditions present -->}
};

% Preload Scanner Thread
\node[procbox=accentB, blade=accentB, text width=8.5cm] (scanL) at (11, 14.5) {
\textbf{\color{accentB}\faSearch\ Preload Scanner Thread}\\
\vspace{4pt}
Scans the raw HTML byte stream solely for \texttt{src} strings. No viewport variables are queried.
};

% Unconditional Push
\node[procbox=accentR, blade=accentR, text width=8.5cm] (pushL) at (11, 10.0) {
\textbf{\color{accentR}\faBolt\ Unconditional Push}\\
\vspace{4pt}
Asset URL is instantly pushed to the Network Thread queue before DOM is even constructed.
};

% Network Thread
\node[procbox=gray, blade=black!60, text width=8.5cm] (netL) at (11, 6.0) {
\textbf{\faGlobe\ Network Thread}\\
\vspace{4pt}
\texttt{GET /wide.jpg}\\
\textit{(Maximum payload dispatched blindly)}
};

% Left Vertical Data Pipes
\draw[datapipe=accentB] (codeL.south) -- (scanL.north);
\draw[datapipe=accentR] (scanL.south) -- (pushL.north);
\draw[datapipe=black!60] (pushL.south) -- (netL.north);

% Left Analytical Annotation (The Baseline Flaw)
\node[bubble=accentR, blade=accentR, text width=6.5cm] (flawL) at (0.9149,10.036) {
\textbf{\color{accentR}The Baseline Flaw:}\\
On mobile endpoints, this algorithm forces the download of a massive \texttt{wide.jpg} file, consuming bandwidth for pixels that will ultimately be suppressed by CSS later.
};

\draw[dashedarrow=accentR] (flawL.east) -- (pushL.west);


% ---------------- RIGHT COLUMN (<picture> + <source> Pipeline) ----------------
\node[align=center] at (33, 22.5) {
    {\LARGE\bfseries\color{charcoal} \texttt{<picture>} Art Direction Pipeline}\\[2pt]
    {\normalsize\color{black!50} (Conditional Fetch via Severing)}
};

% Code Box
\node[codebox, text width=11cm] (codeR) at (33, 18.5) {
<\textcolor{tag}{picture}>\\
\quad<\textcolor{tag}{source} \textcolor{attr}{media}=\textcolor{str}{"(max-width: 600px)"}\\
\quad\quad\quad\quad\textcolor{attr}{srcset}=\textcolor{str}{"tight-crop.jpg"}>\\
\quad<\textcolor{tag}{img} \textcolor{attr}{src}=\textcolor{str}{"wide.jpg"}>\\
</\textcolor{tag}{picture}>
};

% Preload Scanner Thread
\node[procbox=accentB, blade=accentB, text width=9.5cm] (scanR) at (33, 14.5) {
\textbf{\color{accentB}\faSearch\ Preload Scanner Thread}\\
\vspace{4pt}
Identifies \texttt{<source>} tag and extracts the \texttt{media} query string.
};

% HALT Barrier Badge
\node[draw=accentR, fill=accentR!10, text=accentR, font=\scriptsize\bfseries, rounded corners=3pt, line width=1.1pt, inner sep=4pt, blur shadow={shadow opacity=15, shadow yshift=-1pt, shadow blur radius=3pt}] (halt) at (33, 12.4) {\faBan\ HALT: GREEDY FETCH SEVERED};

% The Mechanical Severing (Circuit Breaker Host)
\node[severbox, blade=accentP, text width=10cm] (sevR) at (33, 10.0) {
\textbf{\color{accentP}\faCut\ The Mechanical Severing}\\
\vspace{4pt}
The scanner evaluates the \texttt{media} query instantly against the browser's known environment viewport. It mechanically severs the greedy fetch queue, bypassing any dependency on the DOM or CSSOM.
};

% Viewport Engine Evaluator
\node[procbox=accentG, blade=accentG, text width=9.5cm] (evalR) at (33, 6.0) {
\textbf{\color{accentG}\faMicrochip\ Viewport Engine Evaluator}\\
\vspace{4pt}
Is internal viewport width $\le$ 600px?\\
\textit{(Evaluated per Environment Viewport Metrics)}
};

% Right Column Flow & Circuit Breaker Mechanical Severing
\draw[datapipe=accentB] (codeR.south) -- (scanR.north);
\draw[datapipe=accentB] (scanR.south) -- (halt.north);
\draw[datapipe=accentP, dashed] (halt.south) -- (sevR.north);
\draw[datapipe=accentG] (sevR.south) -- (evalR.north);

% Right Analytical Annotation (The Mechanical Impact)
\node[bubble=accentP, blade=accentP, text width=6.5cm] (impactR) at (43.2925,9.964) {
\textbf{\color{accentP}The Mechanical Impact:}\\
The \texttt{media} attribute acts as a logical circuit breaker. It mathematically ties the network request to the physical endpoint parameters, successfully severing the standard unconditional HTML parsing algorithm.
};

\draw[dashedarrow=accentP] (impactR.west) -- (sevR.east);


% ---------------- NETWORK OUTCOMES (Orthogonal Boolean Branching) ----------------
\node[procbox=gray, blade=accentG, text width=5cm] (netMobile) at (27, 2.0) {
\textbf{\faGlobe\ Network}\\
\vspace{2pt}
\texttt{GET /tight-crop.jpg}
};

\node[procbox=gray, blade=black!50, text width=5cm] (netDesktop) at (39, 2.0) {
\textbf{\faGlobe\ Network}\\
\vspace{2pt}
\texttt{GET /wide.jpg}
};

% Orthogonal Branching Routes (Separated Distinct Boolean Paths)
\draw[datapipe=accentG, rounded corners=6pt] ([xshift=-15pt]evalR.south) -- ([xshift=-15pt]33, 4.0) -- (27, 4.0) -- (netMobile.north);
\node[above, font=\small\bfseries, color=accentG, yshift=2pt] at (30, 4.0) {TRUE (Mobile)};

\draw[datapipe=black!50, rounded corners=6pt] ([xshift=15pt]evalR.south) -- ([xshift=15pt]33, 4.0) -- (39, 4.0) -- (netDesktop.north);
\node[above, font=\small\bfseries, color=black!60, yshift=2pt] at (36, 4.0) {FALSE (Desktop)};

\end{tikzpicture}
\end{document}
